\clearpage
\section*{September 8, 2026: Can ACS earnings explain geographic bunching?}
\textit{Agent-drafted feasibility audit.}

Construction-industry earnings are lower among Bronx and Brooklyn residents
than among Manhattan residents in the pre-policy ACS. The difference remains
when the sample is restricted to full-time, year-round workers. This fits the
hypothesis that a common construction wage floor could impose a larger cost
increase where baseline compensation was lower. It does not establish that
labor on a Bronx or Brooklyn construction site actually cost less.

The distinction matters because the ACS measures residents' annual earnings,
including net self-employment earnings, across all occupations in construction.
Workers can commute to other boroughs, and annual earnings also reflect hours,
weeks, and occupation. Even the full-time, year-round table is not a measure of
hourly pay. The relevant economic variable would be the gap between baseline
job-site compensation and the applicable wage requirement. At the 100-unit
threshold the construction floor is citywide; the higher requirements for
150-plus-unit sites also depend on designated zones.

The audit downloaded the 2019--2023 five-year ACS estimates and their margins
of error for every NYC tract returned by the API and for the five counties.
This window precedes 485-x's enactment, although it includes the pandemic and
extends beyond the project's 2019--2022 filing comparison. The figure uses
directly published county medians, not averages of tract medians. Tract
estimates are often unavailable, especially in Manhattan, and many available
estimates have large margins of error. The saved data preserve these limitations
and distinguish unavailable values from medians reported only above or below
a bound. No annual median is converted into an observed hourly wage.

The geographic bunching result also needs careful wording. Manhattan has a
large spike at 99 in individual filings; it has a much smaller spike in linked
parent totals. A proposed wage mechanism may therefore need to explain project
organization as well as unit choice. Assigning a residential-tract earnings
median to a nearby development would not identify that mechanism.

My recommendation is to retain the ACS results as a descriptive comparison
and investigate employer-location construction payroll using QCEW before
interpreting a tract-level earnings gradient. Public ACS/IPUMS microdata allow
more detailed worker restrictions but identify PUMAs rather than tracts.
Neither alternative has been downloaded or adopted as the project's wage
measure.

{\small Reproduction: run \texttt{make} in
\texttt{tasks/audits/audit\_acs\_construction\_wages/code/}.
The sibling acquisition task owns the Census responses. Tables B24031 and
B24041 report annual earnings in 2023 dollars. The feasibility report contains
coverage counts and the full numerical comparison; standard reports establish
the first saved-data baseline. All cleaned estimates and annotations were
independently reconciled to the responses. The existing home
\texttt{.Renviron} supplies the Census key without storing it in the project;
the separate IPUMS key was not used. Working revision: \texttt{7bbf6d5} plus
this audit. The main analysis and framework are unchanged.}

\clearpage
\begin{center}
\includegraphics[width=\textwidth]{../tasks/audits/audit_acs_construction_wages/output/construction_earnings.pdf}
\vspace{12pt}
\includegraphics[width=\textwidth]{../tasks/audits/audit_acs_construction_wages/output/tract_precision.pdf}
\end{center}
